\chap{Finite observations and realized transformations}\label{sec:observation}
The point-space construction has a converse question. Given a function between coherent realization spaces, can it be represented by compatible maps between finite views? The answer is not automatically yes. This section constructs the relevant category of reindexed maps and identifies its realized functions under an explicit observability hypothesis.

\subsection{An output level may require a finer input level}
Let \(X\) and \(Y\) be sequential set-valued resolutions. Write \(\rho^X_{i,j}:X_j\to X_i\) for composite restriction. A reindexed transport consists of a monotone schedule \(\phi:\N\to\N\) and maps
\[
 f_n:X_{\phi(n)}\longrightarrow Y_n
\]
satisfying, for \(i\le j\),
\begin{equation}\label{eq:scheduled-naturality}
 \rho^Y_{i,j}f_j=f_i\rho^X_{\phi(i),\phi(j)}.
\end{equation}
The schedule records how much input is used for each output level. Monotonicity means that asking for more output never forces the representation to discard input already available. It does not require strict growth, a bounded look-ahead, or an efficiency estimate.

The realized map is
\[
 |f|:\operatorname{Point}(X)\to\operatorname{Point}(Y),\qquad
 |f|(x)_n=f_n(x_{\phi(n)}).
\]
Equation~\eqref{eq:scheduled-naturality} proves that this is a coherent output. A second transport \(g_n:Y_{\psi(n)}\to Z_n\) composes with \(f\) using schedule \(\phi\circ\psi\) and coordinate map \(g_n f_{\psi(n)}\). The order of the schedules is fixed by the types, not by the order in which their letters are written in prose.

\subsection{Equality is witnessed at finite stages}
Two different schedules can represent the same transformation. To compare \((\phi,f)\) and \((\psi,g)\), require that for every output level \(n\) there be a finite level \(k\ge\phi(n),\psi(n)\) with
\begin{equation}\label{eq:eventual-equality}
 f_n\rho^X_{\phi(n),k}=g_n\rho^X_{\psi(n),k}
 \quad\text{as functions on all of }X_k.
\end{equation}
Call this eventual equality. The witnessing \(k\) may depend on \(n\). Equality is required on the whole finite-level fiber, not only on states that happen to extend globally.

\begin{theorem}[The eventual transport category]\label{KS61}
Eventual equality is an equivalence relation and is preserved by composition. Its quotient therefore gives a category of resolutions and eventual transports. Realization descends to a functor into sets, preserving identities and composition.
\end{theorem}
\begin{proof}
Reflexivity and symmetry follow from the defining equation. An equality witnessed at \(k\) persists at every \(l\ge k\), by precomposition with \(\rho^X_{k,l}\). For transitivity, move two witnesses to the common stage given by their maximum and compose the equalities there.

For composition, suppose \(f\sim f'\) and \(g\sim g'\). At output level \(n\), choose a level \(j\) of \(Y\) witnessing \(g\sim g'\). Choose a level \(k\) of \(X\) witnessing \(f\sim f'\) at output \(j\). The naturality equations move the two composite expressions to level \(j\), where the first equality applies, and then to the witness for \(g\sim g'\). This proves equality of the composites at level \(k\). Thus composition is well defined on equivalence classes.

If \(x\) is a coherent input, evaluate \eqref{eq:eventual-equality} at \(x_k\). Restriction replaces its two arguments by \(x_{\phi(n)}\) and \(x_{\psi(n)}\), so eventual equality implies equality of realized outputs. The coordinate formula for composition proves \(|g\circ f|=|g|\circ|f|\), and the identity formula is immediate.
\end{proof}
\source{KS61}
This is an actual quotient category in the formal source. No identification with the full procategory of an arbitrary base category is needed or claimed. The objects here are sequential set-valued resolutions, and the equivalence relation is the finite-stage relation just defined.

\subsection{When global behavior determines the finite presentation}
Say that \(X\) has surjective point projections when
\[
 \forall k\quad \operatorname{ev}_k:\operatorname{Point}(X)\to X_k
 \quad\text{is onto}.
\]
This means that every finite input state really occurs in some coherent global input. It is stronger than nonemptiness of the global space.

\begin{theorem}[Faithfulness on an observable source]\label{KS62}
If \(X\) has surjective point projections, two reindexed maps out of \(X\) are eventually equal exactly when they induce the same realized function. Surjective adjacent bonding maps are a sufficient condition for surjective point projections.
\end{theorem}
\begin{proof}
Only the converse to the previous theorem needs proof. Suppose \(|f|=|g|\), and fix \(n\). Set \(k=\max(\phi(n),\psi(n))\). For any \(a\in X_k\), choose a coherent point \(x\) with \(x_k=a\). Equality of the realized functions at coordinate \(n\) gives
\[
 f_n(x_{\phi(n)})=g_n(x_{\psi(n)}).
\]
Coherence rewrites this as \eqref{eq:eventual-equality} evaluated at \(a\). Since \(a\) was arbitrary, it is the required finite-stage equality.

For the sufficient condition, start with a seed \(a\in X_k\). Its coordinates below \(k\) are determined by restriction. Choose a preimage at level \(k+1\), then a preimage of that choice at level \(k+2\), and continue. The adjacent equations imply every composite restriction equation. This produces a coherent point with the prescribed coordinate. The formal construction uses ordinary classical choice and recursion.
\end{proof}
\source{KS62}
The hypothesis is sufficient, not presented as a characterization of every source on which faithfulness happens to hold. Its role in the proof is precise: it converts a statement about globally realized inputs into a statement about all finite inputs.

Without it, global behavior can lose information. For example, use \(X_n=\{m\in\N:m\ge n\}\) with inclusion bonds and a constant Boolean output tower. The source point space is empty. The two transports that always output false and always output true induce the same function from the empty set, but disagree at every finite stage. Every \(X_k\) is nonempty, so no eventual equality witness exists. This calculation uses the finite-prefix boundary of KS54 and does not identify empty global behavior with equality of finite presentations.

\subsection{The exact image of realization}
For a function \(F:\operatorname{Point}(X)\to\operatorname{Point}(Y)\), define finite observation by the quantifiers
\begin{equation}\label{eq:finite-observation}
 \forall n\ \exists k\ \exists a:X_k\to Y_n\ \forall x:\quad
 F(x)_n=a(x_k).
\end{equation}
The input level is uniform over all global inputs for each fixed output level. It may vary with the output level. This is not the weaker condition that each individual input can be inspected to some input-dependent finite depth.

\begin{theorem}[Characterization of realized transformations]\label{KS63}
If \(X\) has surjective point projections, realization induces a bijection
\[
 \operatorname{EventualHom}(X,Y)
 \simeq
 \{F:\operatorname{Point}(X)\to\operatorname{Point}(Y):F
 \text{ is finitely observed}\}.
\]
\end{theorem}
\begin{proof}
Every reindexed map satisfies \eqref{eq:finite-observation} by using its schedule and coordinate maps. Conversely, choose witnesses \(k_n\) and \(a_n\) for a finitely observed function. The chosen sequence \(k_n\) need not be monotone. Replace it by its running maximum
\[
 \phi(0)=k_0,\qquad \phi(n+1)=\max(\phi(n),k_{n+1}),
\]
and define
\[
 f_n=a_n\rho^X_{k_n,\phi(n)}.
\]
The resulting maps still compute \(F(x)_n\) on coherent inputs.

To prove naturality at \(i\le j\), take an arbitrary state in \(X_{\phi(j)}\) and lift it to a coherent point. Both sides of \eqref{eq:scheduled-naturality} now become the corresponding coordinates of \(F(x)\), which agree by output coherence. Surjectivity of point projections proves the equation on every input state. Thus \((\phi,f)\) is a genuine reindexed map realizing \(F\). The previous theorem identifies two such presentations precisely when their realized functions agree, proving the bijection.
\end{proof}
\source{KS63}
The finite-level compatibility is derived in this proof. It is not inserted as an additional field in the hypothesis that \(F\) is finitely observed. This is one of the useful bridges between a global description of a transformation and an actual presentation in the resolution calculus.

\plate{20_schedule}{Independent finite-observation witnesses need not use a monotone input schedule. Taking the running maximum produces one without changing the realized function. The map at the enlarged level first restricts to the originally sufficient level. This construction gives no efficiency bound on how large those levels are.}{fig:schedule}

\subsection{Binary prefixes: a complete interpretation}
Let \(X_n\) consist of binary words of length \(n\), with prefix restriction. A coherent point is exactly an infinite binary sequence: its length-\(n\) view is its first \(n\) bits. Every prefix extends, for example by appending zeros. Hence the source satisfies the hypothesis of the characterization theorem.

Bitwise complementation is represented at level \(n\) by complementing the first \(n\) bits. The left shift \(F(x)_i=x_{i+1}\) uses the schedule \(\phi(n)=n+1\). A nearest-neighbor difference over \(\mathbb F_2\),
\[
 G(x)_i=x_i+x_{i+1},
\]
also uses \(n+1\) input bits for \(n\) output bits. Applying \(G\) twice uses \(n+2\) bits, and the middle terms cancel:
\[
 G(G(x))_i=x_i+x_{i+2}.
\]
The compatible finite maps make each statement meaningful on the entire coherent sequence, not merely on one chosen truncation. These are worked instances of the scheduled formulas. The finite regression suite checks their prefix compatibility independently.

Not every set function on these points is finitely observed. Consider the Boolean output that is one exactly when every input bit is zero, regarded as a point of the constant Boolean tower. Given any proposed finite level \(k\), compare the all-zero sequence with a sequence whose first one occurs at position \(k\). Their first \(k\) bits agree, but the outputs differ. Therefore no finite witness \(a:X_k\to\{0,1\}\) can compute even that one output coordinate. This rules out a reindexed presentation by the forward direction of the theorem.

The obstruction in this example is one of information, not a runtime benchmark. No amount of computation on the specified finite prefix distinguishes those two inputs. The theorem does not prohibit a different mathematical representation that supplies more information as part of each view. Changing the available view is changing the input model.

\subsection{Cofinal changes also preserve the morphism calculus}
A monotone unbounded schedule \(\theta\) gives the sampled resolution \(X_\theta\). There is a transport from \(X\) to \(X_\theta\) that reads the selected level. For the reverse direction choose, for each \(n\), a level \(j(n)\) with \(n\le\theta(j(n))\), and restrict the available value down to \(X_n\). Taking a least such index gives a monotone schedule.

\begin{theorem}[Cofinal invariance of eventual transports]\label{KS64}
Sampling and recovery are inverse in the eventual transport category. Consequently, for monotone unbounded schedules \(\theta\) and \(\eta\), there is a bijection
\[
 \operatorname{EventualHom}(X,Y)
 \simeq
 \operatorname{EventualHom}(X_\theta,Y_\eta).
\]
On point spaces it agrees with the previously constructed cofinal realization equivalences.
\end{theorem}
\begin{proof}
Each composite samples a finite level and then restricts to a level already required by the original input. At a sufficiently fine common stage, the restriction composition law identifies that composite with the identity. This is exactly eventual equality, so the two transports determine an isomorphism in the quotient category.

Conjugate an arrow \(X\to Y\) by the inverse sampling isomorphism on its source and the sampling isomorphism on its target. Conjugating in the opposite direction gives the inverse map on Hom sets. Realization is a functor, so it sends this conjugation to the corresponding conjugation of functions by the cofinal point equivalences.
\end{proof}
\source{KS64}
This strengthens cofinal invariance from points to a specified category of transformations. It still does not make raw residual norms invariant, identify every global set function with a finite presentation, or preserve the numerical conditioning of a repair. Those are separate questions, addressed by the defect maps in the next sections.

\begin{insight}
For a projection-surjective source, eventual finite presentations correspond exactly to uniformly finite observations at each output level. This is a precise machine-oriented representation theorem. It is not a claim that every represented map is computationally efficient, or that every arbitrary function of an infinite input has such a presentation.
\end{insight}
